% SEGMENT OLP-0623-S01
% Part: history
% Chapter: biographies
% Section: alonzo-church

\documentclass[../../../include/open-logic-section]{subfiles}

\begin{document}

\olfileid{his}{bio}{chu}

\olsection{அலோன்சோ சர்ச்}

\olphoto{church-alonzo}{அலோன்சோ சர்ச்}

அலோன்சோ சர்ச் 1903 ஜூன் 14 அன்று வாஷிங்டன், டி.சி.-யில் பிறந்தார்.
சிறுவயதில் காற்றுத் துப்பாக்கியால் நிகழ்ந்த விபத்தில் ஒரு கண்ணின் பார்வையை
இழந்தார். 1920-இல் கனெக்டிகட்டில் ஆயத்தப் பள்ளிப் படிப்பை முடித்து,
அதே ஆண்டில் பிரின்ஸ்டனில் பல்கலைக்கழகக் கல்வியைத் தொடங்கினார். 1927-இல்
முனைவர் படிப்பை முடித்தார். சில ஆண்டுகள் வெளிநாட்டில் இருந்தபின்
பிரின்ஸ்டனுக்குத் திரும்பினார். மிகுந்த பணிவும் கவனமும் உடையவராகச் சர்ச்
அறியப்பட்டார். கரும்பலகையில் அவர் எழுதுவது மிகவும் நேர்த்தியாக இருக்கும்;
முக்கியமான ஆய்வுக் கட்டுரைகளைப் பாதுகாக்க அவற்றின்மீது டூகோ சிமென்ட்
(ஒளி ஊடுருவும் பசை) கவனமாகப் பூசுவார். கல்விப் பணிக்கு வெளியே
அறிவியல் புனைகதை இதழ்களை வாசிப்பதில் மகிழ்ந்தார்; அவற்றின் எழுத்தில்
பிழை கண்டால் ஆசிரியர்களுக்குக் கடிதம் எழுதவும் தயங்கமாட்டார்.

% SEGMENT OLP-0623-S02
சர்ச்சின் கல்விசார் சாதனைகள் பெருமளவானவை. தமது மாணவர்களான ஸ்டீபன் கிளீனி,
பார்க்லி ரோசர் ஆகியோருடன் சேர்ந்து, பயனுறு கணிக்கத்தன்மைக்கான லாம்டா
கலனம் என்னும் கோட்பாட்டை உருவாக்கினார்; டியூரிங் பொறியை ஆலன் டியூரிங்
உருவாக்கியதிலிருந்து இது தனித்துப் பிறந்தது. கணிக்கத்தன்மையின் இந்த இரண்டு
வரையறைகளும் சமமானவை. இதிலிருந்து இப்போது \emph{சர்ச்--டியூரிங் முன்வைப்பு}
எனப்படும் கருத்து எழுகிறது: இயல் எண்கள் மீதான ஒரு சார்பு பயனுறு முறையில்
கணிக்கத்தக்கது என்றால், அப்போதும் அப்போது மட்டுமே, அதை டியூரிங் பொறி
(அல்லது லாம்டா கலனம்) மூலம் கணிக்க முடியும். மேலும் \emph{சர்ச்சின் தேற்றம்}
எனப்படும் முடிவையும் அவர் நிறுவினார்: முதல்தர வாய்பாடுகளின் செல்லுபடித்
தன்மைக்கான தீர்மானச் சிக்கல் தீர்க்கவியலாதது.

% SEGMENT OLP-0623-S03
சர்ச் முதுமையிலும் தமது பணியைத் தொடர்ந்தார். 1967-இல் பிரின்ஸ்டனை விட்டு
யு.சி.எல்.ஏ.-வுக்குச் சென்றார்; 1990-இல் ஓய்வு பெறும்வரை அங்கே
பேராசிரியராக இருந்தார். 1995 ஆகஸ்ட் 1 அன்று 92 வயதில் மறைந்தார்.

% SEGMENT OLP-0623-S04
\begin{reading}
சர்ச்சின் சுருக்கமான வாழ்க்கை வரலாற்றுக்கு \citet{EndertonND}-ஐப் பார்க்கவும்.
லாம்டா கலனம், தீர்மானச் சிக்கல் (சர்ச்சின் முன்வைப்பு) பற்றிய சர்ச்சின்
மூல எழுத்துகள் \citet{Church1936,Church1936a} ஆகும். 1930-களில்
பிரின்ஸ்டனின் கணிதச் சூழல் குறித்து சர்ச்சுடன் நடத்தப்பட்ட நேர்காணலை
\citet{Aspray1984} பதிவுசெய்கிறது. 1936 முதல் 1979 வரை
\emph{Journal of Symbolic Logic} இதழுக்காகச் சர்ச் தொடர் நூல் மதிப்புரைகள்
எழுதினார். அவை அனைத்தும் ஜான் மேக்ஃபார்லேனின் வலைத்தளத்தில்
காப்பகப்படுத்தப்பட்டுள்ளன \citep{MacFarlane2015}.
\end{reading}

\end{document}
